\section{Related Work}
\label{sec:related}

\begin{table}[t]
\centering
\caption{\textbf{Comparison of research-agent evaluation designs by what each can observe.} Columns: \textbf{Unit scored}—whether scoring is anchored at the final answer (endpoint) or reads the trajectory; \textbf{Artifact-aware}—whether evaluation inspects the agent's generated code, data, and run logs, not its conversational trace or final report alone; \textbf{No reference needed}—whether a task can be evaluated without a known correct answer; Stage attribution—whether a detected failure can be localized to a research stage. Green \yes{} full support, yellow \pt{} partial, red \no{} none.}
\label{tab:comparison}
\setlength{\tabcolsep}{5pt}
\small
\resizebox{\textwidth}{!}{%
\begin{tabular}{llcccccc}
\toprule
\textbf{Benchmark} & \textbf{Domains}
  & \makecell{\textbf{Lifecycle}\\\textbf{coverage}}
  & \makecell{\textbf{Unit}\\\textbf{scored}}
  & \makecell{\textbf{Artifact-}\\\textbf{aware}}
  & \makecell{\textbf{No reference}\\\textbf{needed}}
  & \makecell{\textbf{Stage}\\\textbf{attribution}}
  & \makecell{\textbf{Scoring}\\\textbf{anchor}} \\
\midrule
\rowcolor{sectgray}\multicolumn{8}{l}{\textit{\textbf{Scientific QA \& Knowledge Retrieval}}} \\
GPQA~\cite{rein2023gpqa} / HLE~\cite{phan2025humanity}
  & Multi-Domain  & --          & Endpoint & \no & \no & \no & Exact Match / MCQ \\
SciCode~\cite{tian2024scicode}
  & 6 Domains     & Exec.       & Endpoint & \pt & \no & \no & Target Match \\
\rowcolor{sectgray}\multicolumn{8}{l}{\textit{\textbf{Scientific Paper Reproduction}}} \\
CORE-Bench~\cite{siegel2024core}
  & Science       & Exec.       & Endpoint & \pt & \no & \no & Output Match \\
REPRO-Bench~\cite{hu2025repro}
  & Social Sci.   & Exec.       & Endpoint & \pt & \no & \no & Expert Assert. \\
PaperBench~\cite{starace2025paperbench}
  & 1 (ML)        & Exec.       & \pt\ Rubric & \yes & \no & \pt & Author Rubric \\
\rowcolor{sectgray}\multicolumn{8}{l}{\textit{\textbf{Task Performance \& Engineering Optimization}}} \\
MLE-bench~\cite{chan2025mle}
  & 1 (ML)        & Exec.       & Endpoint & \no & \no & \no & Kaggle \\
NatureBench~\cite{wang2026naturebench}
  & Multi-Science & Exec.+Anal. & Endpoint & \no & \no & \no & Published SOTA \\
ResearchClawBench~\cite{xu2026researchclawbench}
  & 10 Domains    & Exec.+Anal. & Endpoint & \pt & \no & \no & Exec. Tests \\
\rowcolor{sectgray}\multicolumn{8}{l}{\textit{\textbf{Failure Taxonomies \& Annotated Traces}}} \\
MAST~\cite{cemri2026multi}
  & Coding/Math   & Interaction & Trajectory & \no & \yes & \pt & Failure-mode taxonomy \\
TRAIL~\cite{deshpande2025trailtracereasoningagentic}
  & Agentic QA    & Interaction & Trajectory & \no & \yes & \pt & Annotated error spans \\
Who\&When~\cite{zhang2025agent}
  & Multi-Agent   & Interaction & Trajectory & \no & \no & \pt & Decisive-step label \\
\midrule
\rowcolor{rowhl}\textbf{\bench (Ours)} & \textbf{7}
  & \textbf{All 6} & \textbf{Trajectory} & \yes & \yes & \yes
  & \makecell{\textbf{Artifact-aware judge}\\\textbf{(+ metric where available)}} \\
\bottomrule
\end{tabular}
}
\end{table}

\paragraph{Tasks are narrowly-scoped from multiple perspectives.}
Existing discovery benchmarks fall into two groups, and neither captures the
full research lifecycle in real scientific domains. The first places agents in
simulated or fictional worlds, adopted because real experiments are
expensive---but this comes at the cost of
realism~\citep{jansen2024discoveryworldvirtualenvironmentdeveloping,
wang2022scienceworldagentsmarter5th, majumder2024discoverybenchdatadrivendiscoverylarge,
li2026pdagentbenchcharacterizinggroundingarchitecting}. The second stays
grounded in real science but restricts itself to verifiable settings---most
commonly machine learning or software engineering tasks whose outcomes can be
measured against a well-defined target such as human SOTA or an established
metric~\citep{chan2025mle, nathani2025mlgym, starace2025paperbench,
chen2025scienceagentbenchrigorousassessmentlanguage, tian2024scicode,
wang2026naturebench, xu2026researchclawbench}. This prerequisite silently
selects which science gets asked: benchmarks concentrate where a computable
reference is cheap, and coverage of the science domains is correspondingly
narrow. Therefore, the type of task is limited and genuine open-ended discovery tasks are excluded~\citep{wang2026selfrevisingdiscoverysystemsscience}.
Expert-graded case studies escape both
problems~\citep{kirgis2026shadow, li2026scicomp, rawat2026plausible,
horstmann2026neuro, eulig2026cawm}, trading scale for scrutiny---but a few
trajectories in a single domain establish that a failure mode occurs, not how
often, in which systems, or whether better engineering would remove it.
\bench spans seven real-world scientific domains with a mix of
target-anchored and open-ended tasks, covering all six stages of the research
lifecycle.

\paragraph{Endpoint scoring measures performance, not process.}
Nearly all existing benchmarks score an agent by whether its final output
matches a reference: an exact-match label~\citep{rein2023gpqa,
phan2025humanity}, a reproduced result~\citep{siegel2024core, hu2025repro,
starace2025paperbench}, or a published SOTA~\citep{chan2025mle,
wang2026naturebench}. This design has two consequences. First, it invites
reward hacking: circular validation, grader-fitting, and leakage can move the
number without doing the science, and the score cannot distinguish a sound
trajectory from a gamed one~\citep{chan2025mle, hu2025repro,
wang2026naturebench}. Second, the score compresses a long-horizon trajectory
into a single scalar, reporting \emph{that} a run failed but not \emph{why},
\emph{where}, or \emph{how}---so endpoint evaluation can rank systems but
cannot diagnose them. Table~\ref{tab:comparison} contrasts what each evaluation
design can observe: \bench retains metric-based scoring where a target
exists and adds artifact-aware, reference-free annotation precisely where it
does not.

\paragraph{Prior failure diagnoses lack systematic coverage or artifact-level visibility.}
Existing diagnostic work on agent failure falls into three groups, each with a
gap. Expert-graded case studies scrutinize individual trajectories in
depth~\citep{kirgis2026shadow, li2026scicomp, rawat2026plausible,
horstmann2026neuro, eulig2026cawm, trehan2026llmsarentscientistsyet, zhang2026far}, but a few trajectories in a single domain could not lead to generalized failure pattern diagnoses. Although some categorization techniques aimed at scientific discovery offer broader categories, they are frequently derived from small
corpora~\citep{luo2025automateseehiddenpitfalls, eulig2026cawm, bisht2026agenticaiscientistsbuilt} and are not systematic. And the recent line of trace-level
analysis---MAST~\citep{cemri2026multi}, TRAIL~\citep{deshpande2025trailtracereasoningagentic},
Who\&When~\citep{zhang2025agent} and
follow-ups~\citep{rafi2026falat, zhu2026stepfinder},
\citep{zhan2026your}---is the most systematic, but its shared design choices
limit what it can see. These works annotate conversational or tool-call traces,
so failures that are not trace-level---such as computation run on a
circular substrate---are invisible. Because they do not target the autoresearch domain, their stages are interaction phases of a
multi-agent protocol, not steps of the scientific method, so a failure cannot
be attributed to hypothesis formation or experimental design. \taxonomy addresses all three: it is derived from 800 trajectories across eight
systems, annotates the full artifact set---run logs, generated data, code, and
reports---and its stage axis is the research lifecycle itself.